% -----------------------------------------------------------------------------------------------------
\section{小 $x$：Wilson 线（高能）因子化的地位}\label{sec:front-smallx}

小 $x$ 的“因子化”是与共线因子化不同的展开：Regge 极限 $s\gg|t|$（DIS 中 $Q^2$ 固定而 $x\to0$），由 Wilson 线的快度因子化（激波图像）组织——入射体的冲击因子 $\otimes$ 在快度上演化（BK/JIMWLK）的靶 Wilson 线关联函数。它的结果是在程函（Regge）展开之内、固定对数精度上的全阶陈述，而不是有幂次计数控制的领头幂次定理；下面的分级都是相对于这一框架的。领头阶 BK 方程是\cref{thm:smallx-bk}，此处不重述。

\begin{proposition}[NLO BK/JIMWLK 与三圈 BK 核 \tagA]\label{prop:front-nlo-bk}
在程函（Regge）展开的次领头对数精度上：(a) 色偶极子的 NLO 演化含二次与三次非线性；NLO 核是跑动耦合部分（$\propto\beta_0$）、共形部分（$1\to3$ 偶极子跃迁）与一个非共形双对数项之和；其前向线性化与 NLO BFKL 核一致\cite{Balitsky:2007feb}——该文第~VI~节与结论（第~VIII~节）报告的 $\alpha_s^2\zeta(3)\times$（偶极子）的差异，后被同一作者追溯到一个被假定为零的积分的错误求值，修正后消失\cite[第~6~节，“In QCD”]{Balitsky:2009xg}；夸克部分与跑动耦合见\cite{Balitsky:2006wa,Kovchegov:2006vj}；(b) NLO JIMWLK\cite{Balitsky:2013fea,Kovner:2013ona}，其共形对称性\cite{Kovner:2014xia}与“unabridged”形式\cite{Kovner:2014lca}（据摘要）；(c) 平面 $\mathcal N=4$ 超对称 Yang--Mills 中的三圈 BK 核，借助与非全局对数的等价、一种 Lorentz 协变的红外／共线减除方法得到，并与 Pomeron 轨迹和扭度 2 谱核对（\cite{Caron-Huot:2016tzz}，据摘要）。（本命题是关于 Wilson 线算符的演化方程。）
\uses{thm:smallx-bk, thm:smallx-bfkl}
\end{proposition}

\begin{proofsketch}
在激波背景中计算 Wilson 线偶极子的单圈修正（胶子与夸克圈、两个胶子的发射），分离跑动耦合与共形部分，再取线性化前向极限与 NLO BFKL 比较。完整论证见\cite{Balitsky:2007feb,Balitsky:2009xg}。
\end{proofsketch}

\paragraph{地位与缺口。}[A]，在程函框架之内。本书不把它当作与 NLO BFKL 的分歧：\cite{Balitsky:2007feb}中的 $\zeta(3)$ 项已被作者撤回\cite[第~6~节]{Balitsky:2009xg}，也不应归因于纵向截断。

\begin{proposition}[NLO BK 的不稳定性与共线重求和 \tagP]\label{prop:front-bk-collinear}
对现实的初始条件，NLO BK 方程的解不稳定\cite{Lappi:2015fma}；运动学约束与共线重求和可以消除这一不稳定性\cite{Beuf:2014uia,Iancu:2015vea}（均据摘要）。
\uses{prop:front-nlo-bk}
\end{proposition}

\paragraph{地位与缺口。}[P]：各种重求和处方在受控对数之外互不相同；与方案无关的共线改进仍是开放问题（\cref{rem:front-open}）。

\begin{proposition}[{\small 高能极限中的 Wilson 线方法、胶子 Regge 化与 Regge 极限的红外奇异部分} \tagA]\label{prop:front-regge}\hfill\break
(a) 带快度依赖的 Wilson 线集合的程函近似（Balitsky--JIMWLK）再现胶子 Regge 化与 BFKL；从四圈起，BFKL 动力学与红外发散的偶极子公式不相容；在严格平面极限下，理论可以用有限个、依赖过程的 Wilson 线表述而不提及 Regge 化胶子，并在 $\mathcal N=4$ 中精确核对（\cite{Caron-Huot:2013fea}，据摘要，第~4~节已核对）。(b) 高能极限下 $2\to2$ 振幅的红外奇异部分在次领头对数精度上有全阶闭式（双 Reggeon 贡献），预言软反常维度在任意圈数的 Regge 极限（\cite{Caron-Huot:2017zfo}，据摘要）。（本命题只涉及部分子振幅。）
\uses{thm:smallx-bfkl, thm:sudakov-ir-poles}
\end{proposition}

\begin{proofsketch}
把振幅写成快度分级的 Wilson 线关联函数，用 Balitsky--JIMWLK 演化逐阶生成高能对数，并与红外因子化（\cref{thm:sudakov-ir-poles}）的结构比较。完整论证见所引文献。
\end{proofsketch}

\paragraph{地位与缺口。}[A]：程函（Regge）展开之内、固定对数精度上的全阶论证；(a) 只在平面 $\mathcal N=4$ 中精确核对，(b) 在次领头对数精度上对全阶成立；部分来源只读到摘要层次。这不是 QCD 的领头幂次定理：“BK/JIMWLK 是 QCD 的领头幂次定理”是过度表述。

\begin{proposition}[简单 Regge 极点因子化在 NNLL 失效 \tagV]\label{prop:front-regge-pole}
利用红外因子化，$2\to2$ 振幅最简单的 Regge 极点因子化在 $\ln(s/|t|)$ 的次次领头对数（NNLL）精度失效：出现两圈的非对数双极点与三圈的不因子化单对数奇异项；直到次领头对数，硬函数在 $d=4$ 中对全阶为零（\cite{DelDuca:2013ara,DelDuca:2014cya}，据该文摘要）。
\uses{def:smallx-regge, thm:sudakov-ir-poles}
\end{proposition}

\paragraph{地位与缺口。}对超出次领头对数的 Regge 极点因子化为 [V]，证据：基于红外因子化（\cref{thm:sudakov-ir-poles}，[A]）的显式计算；修正后的多 Reggeon 图像是 [A]（\cref{prop:front-regge}(b)）。

\begin{proposition}[小 $x$ 的 TMD 与 Sudakov 对数 \tagA]\label{prop:front-smallx-tmd}
在色玻璃凝聚框架的领头阶与领头对数精度上（均据摘要）：(a) 稀—密碰撞中双喷注的有效 $k_t$ 因子化；Weizsäcker--Williams 与偶极子胶子分布是普适的小 $x$ TMD，在小 $x$ 带依赖过程的规范链\cite{Dominguez:2011wm}；(b) 在饱和形式中把 Sudakov 双对数与小 $x$ 重求和结合\cite{Mueller:2012uf}。(c)（部分结果，\tagP）胶子 TMD 从小 $x$ 到中等 $x$ 的快度演化\cite{Balitsky:2015qba}；小 $x$ 的规范不变 TMD 因子化、只按快度的因子化与 Sudakov--BFKL 插值\cite{Balitsky:2020jzt,Balitsky:2021fer,Balitsky:2023hmh,Balitsky:2026lop}。（高能展开框架内的陈述；靶的非微扰信息作为初始条件输入。）
\uses{thm:smallx-bk}
\end{proposition}

\begin{proofsketch}
在激波背景中计算双喷注截面并在相关极限下与 TMD 形式匹配，识别出两种胶子分布及其规范链；Sudakov 对数由软胶子在饱和背景中的再求和得到。完整论证见所引文献。
\end{proofsketch}

\paragraph{地位与缺口。}(a)(b) 为 [A]（领头阶／领头对数），(c) 为 [P]。

\begin{proposition}[小 $x$ 螺旋度演化：截距与 BER 的分歧 \tagA]\label{prop:front-smallx-helicity}
在双对数近似与大 $N_c$（或大 $N_c\&N_f$）下（据该文摘要），小 $x$ 的螺旋度演化（双对数 $\alpha_s\ln^2(1/x)$）由极化偶极子振幅即次程函 Wilson 线算符给出\cite{Kovchegov:2015pbl,Kovchegov:2016weo}（前者连同勘误），其修正形式见\cite{Cougoulic:2022gbk}；大 $N_c$ 下的精确解给出 $\Delta\Sigma\sim\Delta G\sim g_1\sim(1/x)^{\alpha_h}$，
\begin{equation}\label{eq:front-helicity}
\alpha_h\simeq3.66074\sqrt{\frac{\alpha_sN_c}{2\pi}}\quad\text{\cite{Borden:2023ugd}},\qquad
\alpha_h^{\rm BER}\simeq3.66394\sqrt{\frac{\alpha_sN_c}{2\pi}}\quad\text{（红外演化方程\cite{Bartels:1996wc}）}.
\end{equation}
所求和的 $\Delta\gamma_{GG}(\omega)$ 与固定阶反常维度一致到三圈，从四圈起与 Bartels--Ermolaev--Ryskin（BER）\cite{Bartels:1995iu,Bartels:1996wc}的结果不同；大 $N_c\&N_f$ 见\cite{Adamiak:2023okq}、\cite{Borden:2024bxa}（夸克$\leftrightarrow$胶子跃迁算符）与\cite{Borden:2025ehe}（解析解；与 BER 及 Blümlein--Vogt\cite{Blumlein:1996hb}的细微分歧仍在）。（高能展开框架内的陈述。）
% VF-FRO-17 (formulas typeset with \cite, not \texttt)
\uses{thm:smallx-bk, def:pdf-polarized}
\end{proposition}

\begin{proofsketch}
写出极化偶极子振幅的演化方程（大 $N_c$ 下闭合），在双对数近似下精确求解，读出 Mellin 空间的 $\Delta\gamma_{GG}$ 与其最右奇点。按\cite{Borden:2023ugd}的记法，$\bar a=\alpha_sN_c/(2\pi)$，
\begin{equation}\label{eq:front-dgamma}
\Delta\gamma_{GG}=\frac\omega2\Bigl[1-\sqrt{1-\tfrac{16\bar a}{\omega^2}\sqrt{1-\tfrac{4\bar a}{\omega^2}}}\,\Bigr],\qquad
\Delta\gamma_{GG}^{\rm BER}=\frac\omega2\Bigl[1-\sqrt{1-\tfrac{16\bar a}{\omega^2}\tfrac{1-3\bar a/\omega^2}{1-\bar a/\omega^2}}\,\Bigr]
\end{equation}
（后者为该文转写的 BER 纯胶子大 $N_c$ 形式），展开为 $\Delta\gamma_{GG}=\frac{4\bar a}\omega+\frac{8\bar a^2}{\omega^3}+\frac{56\bar a^3}{\omega^5}+\frac{496\bar a^4}{\omega^7}+\cdots$，BER 的第四项系数为 $504$。截距是外层根号的零点：令 $u=\omega^2/\bar a$，分别得 $u^3-256u+1024=0$ 与 $u^2-17u+48=0$，最大根给出 $\sqrt u=3.66074$ 与 $\sqrt{(17+\sqrt{97})/2}=3.66394$（BER 原文只给出 $3.66$）。完整论证见\cite{Borden:2023ugd,Borden:2025ehe}。
% check: writing/work/W27/s_helicity.py (coefficients 4,8,56,496,4824 vs 4,8,56,504,4920; intercepts)
\end{proofsketch}

\paragraph{地位与缺口。}[A]（双对数近似，大 $N_c$ 或大 $N_c\&N_f$）；与 BER 的分歧是开放问题，本书不作裁决（2026 年的现状见\cite{Borden:2025ehe}）。相关工作（据摘要）：小 $x$ 的胶子轨道角动量\cite{Hatta:2016aoc}与胶子 GPD $E_g$ 的小 $x$ 演化\cite{Hatta:2022bxn}，在相应近似下属全阶论证或部分结果。
